\documentclass[10pt,a4paper]{amsart}
\usepackage{amsmath,amssymb,amsthm,graphicx}
% Retained minimal preamble for the Stacks Project chapter varieties.tex.
% Verbatim excerpt of the chapter preamble: Xy-pic setup (native preamble lines 14-18), the native theorem
% declarations with the mathematical macros (native preamble lines 154-217), and the project own no-op
% definitions of the reference/slogan/history display environments (native wrapper.tex lines 1-3) that the
% counted statement uses for its slogan. Copied mechanically; nothing added, removed or redefined.
% For commutative diagrams we use Xy-pic
\usepackage[all]{xy}

% We use 2cell for 2-commutative diagrams.
\xyoption{2cell}
% Theorem environments.
%
\theoremstyle{plain}
\newtheorem{theorem}[subsection]{Theorem}
\newtheorem{proposition}[subsection]{Proposition}
\newtheorem{lemma}[subsection]{Lemma}

\theoremstyle{definition}
\newtheorem{definition}[subsection]{Definition}
\newtheorem{example}[subsection]{Example}
\newtheorem{exercise}[subsection]{Exercise}
\newtheorem{situation}[subsection]{Situation}

\theoremstyle{remark}
\newtheorem{remark}[subsection]{Remark}
\newtheorem{remarks}[subsection]{Remarks}

\numberwithin{equation}{subsection}

% Macros
%
\def\lim{\mathop{\mathrm{lim}}\nolimits}
\def\colim{\mathop{\mathrm{colim}}\nolimits}
\def\Spec{\mathop{\mathrm{Spec}}}
\def\Hom{\mathop{\mathrm{Hom}}\nolimits}
\def\Ext{\mathop{\mathrm{Ext}}\nolimits}
\def\SheafHom{\mathop{\mathcal{H}\!\mathit{om}}\nolimits}
\def\SheafExt{\mathop{\mathcal{E}\!\mathit{xt}}\nolimits}
\def\Sch{\mathit{Sch}}
\def\Mor{\mathop{\mathrm{Mor}}\nolimits}
\def\Ob{\mathop{\mathrm{Ob}}\nolimits}
\def\Sh{\mathop{\mathit{Sh}}\nolimits}
\def\NL{\mathop{N\!L}\nolimits}
\def\CH{\mathop{\mathrm{CH}}\nolimits}
\def\proetale{{pro\text{-}\acute{e}tale}}
\def\etale{{\acute{e}tale}}
\def\QCoh{\mathit{QCoh}}
\def\Ker{\mathop{\mathrm{Ker}}}
\def\Im{\mathop{\mathrm{Im}}}
\def\Coker{\mathop{\mathrm{Coker}}}
\def\Coim{\mathop{\mathrm{Coim}}}

% Boxtimes
%
\DeclareMathSymbol{\boxtimes}{\mathbin}{AMSa}{"02}

%
% Macros for moduli stacks/spaces
%
\def\QCohstack{\mathcal{QC}\!\mathit{oh}}
\def\Cohstack{\mathcal{C}\!\mathit{oh}}
\def\Spacesstack{\mathcal{S}\!\mathit{paces}}
\def\Quotfunctor{\mathrm{Quot}}
\def\Hilbfunctor{\mathrm{Hilb}}
\def\Curvesstack{\mathcal{C}\!\mathit{urves}}
\def\Polarizedstack{\mathcal{P}\!\mathit{olarized}}
\def\Complexesstack{\mathcal{C}\!\mathit{omplexes}}
% \Pic is the operator that assigns to X its picard group, usage \Pic(X)
% \Picardstack_{X/B} denotes the Picard stack of X over B
% \Picardfunctor_{X/B} denotes the Picard functor of X over B
\def\Pic{\mathop{\mathrm{Pic}}\nolimits}
\def\Picardstack{\mathcal{P}\!\mathit{ic}}
\def\Picardfunctor{\mathrm{Pic}}
\def\Deformationcategory{\mathcal{D}\!\mathit{ef}}
\newenvironment{reference}{}{}
\newenvironment{slogan}{}{}
\newenvironment{history}{}{}
\usepackage{cleveref}
\crefname{originallemma}{original lemma}{original lemmas}
\Crefname{originallemma}{Original lemma}{Original lemmas}
\makeatletter
\newlabel{algebra-lemma-Noetherian-permanence}{{00FN}{source}{}{sourcepointer-1}{}}
\newlabel{algebra-lemma-Noetherian-permanence@cref}{{[originallemma][0][]00FN}{[1][0][]source}{}{}{}}
\makeatother
\crefname{originallemma}{original lemma}{original lemmas}
\Crefname{originallemma}{Original lemma}{Original lemmas}
\makeatletter
\newlabel{algebra-lemma-characterize-dvr}{{00PD}{source}{}{sourcepointer-2}{}}
\newlabel{algebra-lemma-characterize-dvr@cref}{{[originallemma][0][]00PD}{[1][0][]source}{}{}{}}
\makeatother
\crefname{originallemma}{original lemma}{original lemmas}
\Crefname{originallemma}{Original lemma}{Original lemmas}
\makeatletter
\newlabel{algebra-lemma-normal-domain-intersection-localizations-height-1}{{031T}{source}{}{sourcepointer-3}{}}
\newlabel{algebra-lemma-normal-domain-intersection-localizations-height-1@cref}{{[originallemma][0][]031T}{[1][0][]source}{}{}{}}
\makeatother
\crefname{originallemma}{original lemma}{original lemmas}
\Crefname{originallemma}{Original lemma}{Original lemmas}
\makeatletter
\newlabel{lemma-normalization-locally-algebraic}{{0BXR}{source}{}{sourcepointer-4}{}}
\newlabel{lemma-normalization-locally-algebraic@cref}{{[originallemma][0][]0BXR}{[1][0][]source}{}{}{}}
\makeatother
\crefname{originallemma}{original lemma}{original lemmas}
\Crefname{originallemma}{Original lemma}{Original lemmas}
\makeatletter
\newlabel{lemma-regular-functions-proper-variety}{{04L2}{source}{}{sourcepointer-5}{}}
\newlabel{lemma-regular-functions-proper-variety@cref}{{[originallemma][0][]04L2}{[1][0][]source}{}{}{}}
\makeatother
\crefname{originallemma}{original lemma}{original lemmas}
\Crefname{originallemma}{Original lemma}{Original lemmas}
\makeatletter
\newlabel{morphisms-lemma-composition-proper}{{01W3}{source}{}{sourcepointer-6}{}}
\newlabel{morphisms-lemma-composition-proper@cref}{{[originallemma][0][]01W3}{[1][0][]source}{}{}{}}
\makeatother
\crefname{originallemma}{original lemma}{original lemmas}
\Crefname{originallemma}{Original lemma}{Original lemmas}
\makeatletter
\newlabel{morphisms-lemma-finite-proper}{{01WN}{source}{}{sourcepointer-7}{}}
\newlabel{morphisms-lemma-finite-proper@cref}{{[originallemma][0][]01WN}{[1][0][]source}{}{}{}}
\makeatother
\crefname{originallemma}{original lemma}{original lemmas}
\Crefname{originallemma}{Original lemma}{Original lemmas}
\makeatletter
\newlabel{morphisms-lemma-finite-type-noetherian}{{01T6}{source}{}{sourcepointer-8}{}}
\newlabel{morphisms-lemma-finite-type-noetherian@cref}{{[originallemma][0][]01T6}{[1][0][]source}{}{}{}}
\makeatother
\crefname{originallemma}{original lemma}{original lemmas}
\Crefname{originallemma}{Original lemma}{Original lemmas}
\makeatletter
\newlabel{morphisms-lemma-locally-finite-type-characterize}{{01T2}{source}{}{sourcepointer-9}{}}
\newlabel{morphisms-lemma-locally-finite-type-characterize@cref}{{[originallemma][0][]01T2}{[1][0][]source}{}{}{}}
\makeatother
\crefname{originaldefinition}{original definition}{original definitions}
\Crefname{originaldefinition}{Original definition}{Original definitions}
\makeatletter
\newlabel{properties-definition-integral}{{01OK}{source}{}{sourcepointer-10}{}}
\newlabel{properties-definition-integral@cref}{{[originaldefinition][0][]01OK}{[1][0][]source}{}{}{}}
\makeatother
\crefname{originallemma}{original lemma}{original lemmas}
\Crefname{originallemma}{Original lemma}{Original lemmas}
\makeatletter
\newlabel{properties-lemma-Noetherian-topology}{{01OZ}{source}{}{sourcepointer-11}{}}
\newlabel{properties-lemma-Noetherian-topology@cref}{{[originallemma][0][]01OZ}{[1][0][]source}{}{}{}}
\makeatother
\crefname{originallemma}{original lemma}{original lemmas}
\Crefname{originallemma}{Original lemma}{Original lemmas}
\makeatletter
\newlabel{properties-lemma-codimension-local-ring}{{02IZ}{source}{}{sourcepointer-12}{}}
\newlabel{properties-lemma-codimension-local-ring@cref}{{[originallemma][0][]02IZ}{[1][0][]source}{}{}{}}
\makeatother
\crefname{originallemma}{original lemma}{original lemmas}
\Crefname{originallemma}{Original lemma}{Original lemmas}
\makeatletter
\newlabel{properties-lemma-locally-Noetherian}{{01OW}{source}{}{sourcepointer-13}{}}
\newlabel{properties-lemma-locally-Noetherian@cref}{{[originallemma][0][]01OW}{[1][0][]source}{}{}{}}
\makeatother
\crefname{originallemma}{original lemma}{original lemmas}
\Crefname{originallemma}{Original lemma}{Original lemmas}
\makeatletter
\newlabel{properties-lemma-locally-normal}{{033J}{source}{}{sourcepointer-14}{}}
\newlabel{properties-lemma-locally-normal@cref}{{[originallemma][0][]033J}{[1][0][]source}{}{}{}}
\makeatother
\crefname{originallemma}{original lemma}{original lemmas}
\Crefname{originallemma}{Original lemma}{Original lemmas}
\makeatletter
\newlabel{topology-lemma-Noetherian}{{0052}{source}{}{sourcepointer-15}{}}
\newlabel{topology-lemma-Noetherian@cref}{{[originallemma][0][]0052}{[1][0][]source}{}{}{}}
\makeatother
\title{Finiteness of the group of units modulo constants for a normal open subscheme of a proper variety}
\author{The Stacks Project Authors}
\date{stacks-project:varieties.tex:local-snapshot-20261011; GNU FDL 1.2 or later}
\begin{document}
\maketitle
\noindent\textit{Editorial note.} All delivered mathematical text is the author's own, copied from the pinned Stacks Project chapter file varieties.tex; no statement or proof step was rewritten, reordered or removed.\par\medskip
\noindent\textit{Editorial note.} The proof contains the parenthetical '(tiny detail omitted)' at the step where the restriction of the normalization over the normal X is called an isomorphism; the parenthetical is reproduced verbatim and the surrounding reduction is complete as written.\par\medskip
\noindent\textit{Editorial note.} The counted spans carry 15 navigation stubs; none of them replaces a load-bearing step. No bibliography citation occurs in the spans.\par\medskip
\noindent\textit{Editorial note.} The preamble is the chapter's retained minimal preamble (sources/stacks-varieties/retained-preamble.tex lines 1-77), a verbatim excerpt of the native preamble.\par\medskip
\noindent\textit{Editorial note.} The label lemma-open-in-normal-proper is defined only in varieties.tex, so no cross-chapter restatement had to be weighed for this card.\par\medskip
\noindent\textit{Allowed standard prerequisites (editorial).}\begin{itemize}
\item Normalization of a variety and finiteness of the normalization morphism. Same chapter, Lemma lemma-normalization-locally-algebraic, invoked at proof lines 5168-5171. Supplies the reduction that lets X-bar be assumed normal.
\item Characterization of discrete valuation rings among one dimensional normal Noetherian local domains. Algebra, Lemma algebra-lemma-characterize-dvr, invoked at proof line 5201. Supplies the valuations v\_j whose product receives O(X)\textasciicircum{}*.
\item Normality as the intersection of the localizations at height one primes. Algebra, Lemma algebra-lemma-normal-domain-intersection-localizations-height-1, invoked at proof lines 5231-5233. Supplies that an element lying in every height one localization of an affine piece lies in the ring.
\item A proper variety over an algebraically closed field has only constants as global functions. Same chapter, Lemma lemma-regular-functions-proper-variety, invoked at proof lines 5236-5238. Supplies that the extended function lies in k, which identifies the kernel with k\textasciicircum{}*.
\end{itemize}\medskip
\ifdefined\hypertarget\hypertarget{sourcepointer-1}{}\fi
\noindent\textit{Source reference.} Original lemma 00FN: sources/stacks-chow/tags.txt, source lines 577--577. Pointer outside the counted span, invoked once for context; the load-bearing argument of the counted statement is the author's proof reproduced in full in the proof span.\par\medskip
\ifdefined\hypertarget\hypertarget{sourcepointer-2}{}\fi
\noindent\textit{Source reference.} Original lemma 00PD: sources/stacks-chow/tags.txt, source lines 927--927. Pointer outside the counted span, invoked once for context; the load-bearing argument of the counted statement is the author's proof reproduced in full in the proof span.\par\medskip
\ifdefined\hypertarget\hypertarget{sourcepointer-3}{}\fi
\noindent\textit{Source reference.} Original lemma 031T: sources/stacks-chow/tags.txt, source lines 3996--3996. Pointer outside the counted span, invoked once for context; the load-bearing argument of the counted statement is the author's proof reproduced in full in the proof span.\par\medskip
\ifdefined\hypertarget\hypertarget{sourcepointer-4}{}\fi
\noindent\textit{Source reference.} Original lemma 0BXR: sources/stacks-varieties/pinned-original-source.tex, source lines 4903--4903. Pointer outside the counted span, invoked once for context; the load-bearing argument of the counted statement is the author's proof reproduced in full in the proof span.\par\medskip
\ifdefined\hypertarget\hypertarget{sourcepointer-5}{}\fi
\noindent\textit{Source reference.} Original lemma 04L2: sources/stacks-varieties/pinned-original-source.tex, source lines 4876--4876. Pointer outside the counted span, invoked once for context; the load-bearing argument of the counted statement is the author's proof reproduced in full in the proof span.\par\medskip
\ifdefined\hypertarget\hypertarget{sourcepointer-6}{}\fi
\noindent\textit{Source reference.} Original lemma 01W3: sources/stacks-chow/tags.txt, source lines 2468--2468. Pointer outside the counted span, invoked once for context; the load-bearing argument of the counted statement is the author's proof reproduced in full in the proof span.\par\medskip
\ifdefined\hypertarget\hypertarget{sourcepointer-7}{}\fi
\noindent\textit{Source reference.} Original lemma 01WN: sources/stacks-chow/tags.txt, source lines 2488--2488. Pointer outside the counted span, invoked once for context; the load-bearing argument of the counted statement is the author's proof reproduced in full in the proof span.\par\medskip
\ifdefined\hypertarget\hypertarget{sourcepointer-8}{}\fi
\noindent\textit{Source reference.} Original lemma 01T6: sources/stacks-chow/tags.txt, source lines 2363--2363. Pointer outside the counted span, invoked once for context; the load-bearing argument of the counted statement is the author's proof reproduced in full in the proof span.\par\medskip
\ifdefined\hypertarget\hypertarget{sourcepointer-9}{}\fi
\noindent\textit{Source reference.} Original lemma 01T2: sources/stacks-chow/tags.txt, source lines 2359--2359. Pointer outside the counted span, invoked once for context; the load-bearing argument of the counted statement is the author's proof reproduced in full in the proof span.\par\medskip
\ifdefined\hypertarget\hypertarget{sourcepointer-10}{}\fi
\noindent\textit{Source reference.} Original definition 01OK: sources/stacks-chow/tags.txt, source lines 2197--2197. Pointer outside the counted span, invoked once for context; the load-bearing argument of the counted statement is the author's proof reproduced in full in the proof span.\par\medskip
\ifdefined\hypertarget\hypertarget{sourcepointer-11}{}\fi
\noindent\textit{Source reference.} Original lemma 01OZ: sources/stacks-chow/tags.txt, source lines 2212--2212. Pointer outside the counted span, invoked once for context; the load-bearing argument of the counted statement is the author's proof reproduced in full in the proof span.\par\medskip
\ifdefined\hypertarget\hypertarget{sourcepointer-12}{}\fi
\noindent\textit{Source reference.} Original lemma 02IZ: sources/stacks-chow/tags.txt, source lines 3318--3318. Pointer outside the counted span, invoked once for context; the load-bearing argument of the counted statement is the author's proof reproduced in full in the proof span.\par\medskip
\ifdefined\hypertarget\hypertarget{sourcepointer-13}{}\fi
\noindent\textit{Source reference.} Original lemma 01OW: sources/stacks-chow/tags.txt, source lines 2209--2209. Pointer outside the counted span, invoked once for context; the load-bearing argument of the counted statement is the author's proof reproduced in full in the proof span.\par\medskip
\ifdefined\hypertarget\hypertarget{sourcepointer-14}{}\fi
\noindent\textit{Source reference.} Original lemma 033J: sources/stacks-chow/tags.txt, source lines 4058--4058. Pointer outside the counted span, invoked once for context; the load-bearing argument of the counted statement is the author's proof reproduced in full in the proof span.\par\medskip
\ifdefined\hypertarget\hypertarget{sourcepointer-15}{}\fi
\noindent\textit{Source reference.} Original lemma 0052: sources/stacks-chow/tags.txt, source lines 194--194. Pointer outside the counted span, invoked once for context; the load-bearing argument of the counted statement is the author's proof reproduced in full in the proof span.\par\medskip
\begin{lemma}
\label{lemma-open-in-normal-proper}
Let $k$ be an algebraically closed field.
Let $\overline{X}$ be a proper variety over $k$.
Let $X \subset \overline{X}$ be an open subscheme.
Assume $X$ is normal.
Then $\mathcal{O}^*(X)/k^*$ is a finitely generated abelian group.
\end{lemma}

\begin{proof}
Since the statement only concerns $X$, we may replace $\overline{X}$
by a different proper variety over $k$. Let $\nu : \overline{X}^\nu \to
\overline{X}$ be the normalization morphism. By
Lemma \ref{lemma-normalization-locally-algebraic}
we have that $\nu$ is finite and $\overline{X}^\nu$ is a
variety. Since $X$ is normal, we see that $\nu^{-1}(X) \to X$ is an
isomorphism (tiny detail omitted). Finally, we see that
$\overline{X}^\nu$ is proper over $k$ as a finite morphism is proper
(Morphisms, Lemma \ref{morphisms-lemma-finite-proper})
and compositions of proper morphisms
are proper (Morphisms, Lemma \ref{morphisms-lemma-composition-proper}).
Thus we may and do assume $\overline{X}$ is normal.

\medskip\noindent
We will use without further mention that for any affine open $U$ of
$\overline{X}$ the ring $\mathcal{O}(U)$ is a finitely generated
$k$-algebra, which is Noetherian, a domain and normal, see
Algebra, Lemma \ref{algebra-lemma-Noetherian-permanence},
Properties, Definition \ref{properties-definition-integral},
Properties, Lemmas \ref{properties-lemma-locally-Noetherian} and
\ref{properties-lemma-locally-normal},
Morphisms, Lemma \ref{morphisms-lemma-locally-finite-type-characterize}.

\medskip\noindent
Let $\xi_1, \ldots, \xi_r$ be the generic points of the complement of $X$
in $\overline{X}$. There are finitely many since $\overline{X}$ has a
Noetherian underlying topological space (see
Morphisms, Lemma \ref{morphisms-lemma-finite-type-noetherian},
Properties, Lemma \ref{properties-lemma-Noetherian-topology}, and
Topology, Lemma \ref{topology-lemma-Noetherian}).
For each $i$ the local ring $\mathcal{O}_i = \mathcal{O}_{X, \xi_i}$
is a normal Noetherian local domain (as a localization of a
Noetherian normal domain). Let $J \subset \{1, \ldots, r\}$ be the set of
indices $i$ such that $\dim(\mathcal{O}_i) = 1$. For $j \in J$ the
local ring $\mathcal{O}_j$ is a discrete valuation ring, see
Algebra, Lemma \ref{algebra-lemma-characterize-dvr}.
Hence we obtain a valuation
$$
v_j : k(\overline{X})^* \longrightarrow \mathbf{Z}
$$
with the property that $v_j(f) \geq 0 \Leftrightarrow f \in \mathcal{O}_j$.

\medskip\noindent
Think of $\mathcal{O}(X)$ as a sub $k$-algebra of $k(X) = k(\overline{X})$.
We claim that the kernel of the map
$$
\mathcal{O}(X)^* \longrightarrow
\prod\nolimits_{j \in J} \mathbf{Z},
\quad
f \longmapsto \prod v_j(f)
$$
is $k^*$. It is clear that this claim proves the lemma.
Namely, suppose that $f \in \mathcal{O}(X)$ is an element of the kernel.
Let $U = \Spec(B) \subset \overline{X}$ be any affine open.
Then $B$ is a Noetherian normal domain.
For every height one prime $\mathfrak q \subset B$ with corresponding
point $\xi \in X$ we see that either $\xi = \xi_j$ for some $j \in J$
or that $\xi \in X$. The reason is that
$\text{codim}(\overline{\{\xi\}}, \overline{X}) = 1$ by
Properties, Lemma \ref{properties-lemma-codimension-local-ring}
and hence if $\xi \in \overline{X} \setminus X$ it must be a
generic point of $\overline{X} \setminus X$, hence equal to some
$\xi_j$, $j \in J$.
We conclude that $f \in \mathcal{O}_{X, \xi} = B_{\mathfrak q}$
in either case as $f$ is in the kernel of the map. Thus
$f \in \bigcap_{\text{ht}(\mathfrak q) = 1} B_{\mathfrak q} = B$, see
Algebra, Lemma
\ref{algebra-lemma-normal-domain-intersection-localizations-height-1}.
In other words, we see that
$f \in \Gamma(\overline{X}, \mathcal{O}_{\overline{X}})$.
But since $k$ is algebraically closed we conclude that
$f \in k$ by
Lemma \ref{lemma-regular-functions-proper-variety}.
\end{proof}

\end{document}
